\section{问题重述}

\subsection{任务概述}
在乡村振兴背景下，需要在 2024--2030 年七年期内，为若干异构耕地（露天/大棚）按季节选择作物并安排种植面积，在满足农艺规则（轮作、防重茬、作物—地块适宜性）与产销约束的前提下，最大化长期经济收益并控制不确定性风险。题目要求分别解决：
- 问题一（确定性最优化）：在确定性参数下获得七年累计净收益最大化的可执行方案；
- 问题二（不确定性与保底收益）：在区间扰动下，通过鲁棒优化获得“最坏情景收益”下界尽可能大的方案；
- 问题三（价格反馈与动态决策）：考虑供给—价格弹性与逐年信息滚动更新，形成风险可控的滚动优化策略，并评估尾部损失（CVaR）。

\subsection{数据与预处理}
根据题面及附件提供的数据，包含：地块编号、面积、季节可用性、大棚标识、作物—地块适宜性；作物单产、价格、成本、销量上限以及（若有）超产价格折减机制或价格区间。以 2023 年为基期。为保证求解与评估稳定性：统一单位口径；识别并修正显著异常价格点；对少量缺失值按分布一致性进行温和填补；保留季节与地块层面的结构特征。

\subsection{问题拆解}
- 确定性（问题一）：构建多周期混合整数规划（Multi-period MIP），以“年–季–地块–作物”面积为连续变量，配合作物选择二元变量控制同季不交叠、重茬与分散度；目标为七年累计净收益最大化。
- 鲁棒（问题二）：考虑单产、价格、成本、销量上限的区间波动，引入盒式/预算不确定集（保守度 \(\Gamma\)）并进行鲁棒对偶重构，使目标成为“最坏情景”净收益下界最大。
- 动态（问题三）：在供给—价格弹性与信息逐年更新背景下，采用“参数更新—情景生成—再优化”的滚动框架，并以 CVaR@95\% 控制尾部风险，比较与问题一/二的收益—风险权衡。

\subsection{模型轮廓与求解思路}
核心约束包括：地块面积与季节可用性、作物—地块适宜性与大棚限制、同季同地块不交叠、禁止重茬与轮作周期、产销匹配与销量上限、（若给定）超产价格折减。采用商用求解器（如 Gurobi）对确定性与鲁棒等价模型求解；动态部分以年度滚动方式迭代，外层蒙特卡洛评估稳健性与分位表现。通过变量预筛、适宜性削减与季节—地块分解启发式降低规模，保证分钟级可解或年度级迭代可用。

\subsection{评价指标与结果呈现}
- 经济性：七年累计/年均收益、收益波动率、CVaR@95\%、最坏情景收益下界；
- 农艺与实施性：重茬/轮作约束满足率、作物/地块切换频率、地块分散度；
- 稳健性：对价格/单产/销量扰动的灵敏度、\(\Gamma\)—收益前沿、蒙特卡洛失效率与分位收益。
结果以“年/季/地块/作物/面积/产量/销量/收入/成本/净收益”等表格与关键曲线呈现。

\subsection{极简符号表}
为避免与“假设与符号”(见第\ref{sec:assumptions}章)及“符号表（扩展）”(见第\ref{sec:assumptions-symbols}章)的详细内容重复，本节仅保留最小必要条目以便快速阅读：

\noindent\textbf{集合（最小）}：\(\mathcal{T}\)、\(\mathcal{S}\)、\(\mathcal{I}\)、\(\mathcal{K}\)；如涉及大棚与轮作，使用 \(\mathcal{G}\)、\(\mathcal{K}_{\mathrm{green}}\)、\(\mathcal{K}_{\mathrm{legume}}\)。

\noindent\textbf{变量（最小）}：面积 \(x_{t s i k}\)、选择 \(z_{t s i k}\)、销量 \(q_{t s k}\)、年净收益 \(r_t\)。

\noindent\textbf{参数（最小）}：面积上限 \(A_{i s}\)、适宜性 \(\eta_{i s k}\)、单产 \(y_{t s k}\)、单价 \(p_{t s k}\)、成本 \(c_{t s k}\)、销量上限 \(u_{t s k}\)。

\noindent 其余（如折价系数 \(\phi_{t s k}\)、轮作间隔 \(L_k\)、鲁棒度 \(\Gamma\)、CVaR 置信水平 \(\alpha\) 等）请见第\ref{sec:assumptions-symbols}节。